\par\clearpage
\SourceLeaf{582}\GaussVSourceFolio{576}{L}{MAGNETISCHE BEOBACHTUNGEN.}
\noindent 1836 beobachtet wäre; sehr wahrscheinlich hätte sich dann über die zwar an den meisten Orten durchscheinende, aber in München auffallend grössere Bewegung um $2^{\mathrm h}10'$ schon mit mehr Sicherheit urtheilen lassen, als jetzt möglich ist.
\par\smallskip\noindent\hspace*{.07\linewidth}\makebox[.93\linewidth][s]{---\hfill---\hfill---\hfill---\hfill---\hfill---\hfill---\hfill---\hfill---\hfill---\hfill---\hfill---\hfill---\hfill---\hfill---}\par\smallskip

\vspace{.075\textheight}
\pdfbookmark[1]{Erläuterungen zu den Terminszeichnungen und den Beobachtungszahlen. 1837.}{erlaeuterungen-termine-1837}
\begin{center}
\resizebox{.845\linewidth}{!}{\textit{Erläuterungen zu den Terminszeichnungen und den Beobachtungszahlen.}}\par
\smallskip[\textit{Im Auszuge.}]
\end{center}
\begin{center}
\vbox{\offinterlineskip
\centerline{\rule{.780\linewidth}{.35pt}}
\kern4pt
\centerline{\resizebox{.680\linewidth}{!}{\small Resultate aus den Beobachtungen des magnetischen Vereins.\quad 1837.\quad VIII.}}
\kern4.5pt
\centerline{\rule{.780\linewidth}{.35pt}}}
\end{center}

Es sind im Jahre 1837 sieben vierundzwanzigstündige Termine abgehalten, da zu den sechs gewöhnlichen noch ein ausserordentlicher am 31. August hinzugekommen ist. Die in den ---\quad---\quad--- Tafeln mitgetheilten Zahlen enthalten 80 Beobachtungsreihen für die Variationen der Declination aus 16 verschiedenen Beobachtungsorten, nemlich Altona, Augsburg, Berlin, Breda, Breslau, Copenhagen, Dublin, Freiberg, Göttingen, Leipzig, Mailand, Marburg, München, Petersburg, Stockholm und Upsala. Es sind uns ausserdem noch einige andere Beobachtungsreihen zugekommen, die wegen zu späten Empfangs nicht mit abgedruckt werden konnten.
\par\smallskip\noindent\hspace*{.07\linewidth}\makebox[.93\linewidth][s]{---\hfill---\hfill---\hfill---\hfill---\hfill---\hfill---\hfill---\hfill---\hfill---\hfill---\hfill---\hfill---\hfill---\hfill---}\par\smallskip
Zur Ansetzung eines ausserordentlichen Termins am 31. August gab Veranlassung die Nachricht, dass Hr.\ Prof.\ \textsc{Parrot} an diesem Tage (so wie an einigen vorhergehenden) die Variation der magnetischen Declination auf dem Nordcap beobachten würde. Die Einladung zur Theilnahme wurde daher, so weit es die Kürze der Zeit verstattete, verbreitet. Es sind dadurch recht interessante Beobachtungsreihen eingebracht, aber die Beobachtungen vom Nordcap selbst sind bisher nicht zu unsrer Kenntniss gekommen.\hfill\mbox{---\enspace---\enspace---\enspace---\enspace---\enspace---\enspace---}

Für den Novembertermin war die sonst befolgte Bestimmung dahin abgeändert, dass er auf den 13. verlegt wurde. Es geschah dies in Folge eines Gesprächs mit Hrn \textsc{von Humboldt} über die Möglichkeit, dass an den Monatstagen, die in mehrern frühern Jahren durch eine ausserordentliche Menge von~Stern-\par
\clearpage
\begingroup\renewcommand{\footnoterule}{\kern-4.35pt\hrule width .253\columnwidth height .35pt\kern4pt}
\SourceLeaf{583}\GaussVSourceFolio{577}{R}{MAGNETISCHE BEOBACHTUNGEN.}
\noindent schnuppen ausgezeichnet gewesen waren, vielleicht auch ungewöhnliche magnetische Bewegungen eintreten könnten. Diese Erwartung hat sich jedoch in sofern nicht bestätigt, als die magnetischen Bewegungen während dieser vierundzwanzig Stunden, wenn gleich sehr beträchtlich, doch nicht grösser als in vielen frühern Terminen zu jeder andern Jahreszeit gewesen sind. Dagegen waren am vorhergehenden und am folgenden Abend an mehrern Orten sehr starke und schnell wechselnde Anomalien in der magnetischen Declination beobachtet, zwischen denen und den Sternschnuppenerscheinungen man aber nicht berechtigt ist, einen Zusammenhang anzunehmen, da jene nur die gewöhnlichen Begleiter von Nordlichtern sind, und sehr glänzende Nordlichter in diesen beiden Nächten wirklich Statt gefunden haben\expandafter\def\csname HistoricalNoteText583\endcsname{Es sind uns die magnetischen Beobachtungen vom 12. November aus Upsala, Leipzig, Breslau und Mailand, und vom 14. November aus Upsala, Dublin, Berlin, Breslau und Mailand mitgetheilt. Ähnlichkeit der Bewegungen ist hier an einigen Stellen unverkennbar, an andern nur schwach durchscheinend. Aber es wiederholt sich hier die auch schon bei anderer Gelegenheit gemachte Bemerkung, dass unter solchen Umständen die Bewegungen viel zu schnell wechseln, als dass Beobachtungen von fünf zu fünf Minuten, oder gar in noch weitern Zwischenzeiten, ein treues Bild davon geben könnten.}\HistoricalNote{583}.
\par\smallskip\noindent\hspace*{.07\linewidth}\makebox[.93\linewidth][s]{---\hfill---\hfill---\hfill---\hfill---\hfill---\hfill---\hfill---\hfill---\hfill---\hfill---\hfill---\hfill---\hfill---\hfill---\hfill---}\par\smallskip
In den Terminen vom Julius, August und November sind in Göttingen nun auch die Variationen der Intensität mit dem Bifilar-Magnetometer vollständig beobachtet. In die Tafeln sind aber nicht die unmittelbar beobachteten Scalentheile selbst aufgenommen, sondern ihre Differenzen von dem grössten in jedem Termine vorgekommenen Werthe. Da in den beiden ersten Terminen diejenige transversale Lage Statt hatte, für welche wachsenden Scalentheilen abnehmende Intensitäten entsprechen, so zeigen hier die Zahlen an, um wie viel die jedesmalige Intensität grösser war, als die kleinste des Termins, und zwar in solchen Einheiten gemessen, wovon für den Juliustermin 22000 auf die kleinste selbst kommen. Da es für jetzt, so lange dergleichen Beobachtungen nur an Einem Orte gemacht werden, auf die schärfste Angabe des \textit{absoluten} Werthes der Scalentheile eben nicht ankommt, so waren zu dem Ende für den Augusttermin keine neuen Bestimmungen gemacht. Vor dem Novembertermine war dies aber geschehen: die geänderte absolute Zahl steht im Zusammenhange mit dem Verluste, welchen der Magnetismus des Stabes in den vier Monaten erlitten hatte. Nur muss bemerkt werden, dass im Novembertermin die Zahlen die Bedeutung haben,\par
\GaussVSourceGathering{73}
\clearpage
\endgroup
\SourceLeaf{584}\GaussVSourceFolio{578}{L}{MAGNETISCHE BEOBACHTUNGEN.}
\noindent um wie viel die jedesmalige Intensität kleiner ist, als die grösste, diese selbst $=18290$ angenommen. Da nemlich in diesem Termine der Stab die entgegengesetzte transversale Lage hatte, für welche Intensität und Scalentheile zugleich wachsen, so hätten, behuf gleichförmiger Bedeutung der Zahlen, die einzelnen unmittelbar beobachteten Scalentheile nicht mit dem Maximum, sondern mit dem Minimum verglichen werden müssen, was durch Versehen nicht beachtet ist.

Auf den Tafeln ---\quad---\quad---\quad--- sind die beobachteten Intensitätsänderungen graphisch dargestellt. Einmal in Einer Curve, unter welcher die Curve für die gleichzeitigen Declinationsänderungen in Göttingen wiederholt ist, wodurch die oben [S.\ 365 d.\ B.] gemachte Bemerkung augenfällig wird, dass nemlich um die Zeit starker Störungen der Declination meistens auch starke Anomalien der Intensität eintreten. Zweitens ist auch der Gang der Veränderungen beider Elemente in Eine Curve zusammengefasst, wodurch man ein anschauliches Bild der Veränderungen des horizontalen Theils der erdmagnetischen Kraft während jedes Termins erhält. Nur haben, um Verwirrung wegen der vielfachen Durchkreuzungen zu vermeiden, die Bewegungen im Julius- und Augusttermin in zwei Stücken, die im Novembertermin in drei Stücken gezeichnet werden müssen, wobei ausserdem zu grösserer Erleichterung der Übersicht jedes Stück zur einen Hälfte in ausgezogenen Linien, zur andern Hälfte punktirt dargestellt ist. Nach dem, was bereits oben [S.\ 366 d.\ B.] bemerkt ist, werden diese Darstellungen einer weitern Erläuterung nicht bedürfen.

Über die Ausbeute selbst können hier nur noch einige Bemerkungen Platz finden. Die ausserordentlich grosse Ähnlichkeit der gleichzeitigen Declinationsbewegungen, an verschiedenen Orten, meistens bis zu den kleinsten Schattirungen herab, bestätigt sich hier wieder eben so schön, wie bei den Beobachtungen des vorhergehenden Jahres. Allein, es werden doch auch hin und wieder schon erhebliche Unterschiede kenntlich, besonders in denjenigen Terminen, wo die Beobachtungen sich über einen noch weitern Umfang erstrecken, obwohl diese Ausdehnung noch immer zu klein, und die Anzahl weit von einander entlegener Örter zu gering erscheint, als dass man schon Schlüsse über die Sitze der Ursachen der einzelnen Bewegungen darauf gründen dürfte. Immerhin würde zwar die nähere Betrachtung mancher einzelnen Bewegungen, zumal von denjenigen Terminen, wo in Göttingen zugleich die Intensitätsänderungen beobachtet sind, zu allerlei Bemerkungen und selbst allgemeinen Betrachtungen Anlass geben~kön-\par
\clearpage
\SourceLeaf{585}\GaussVSourceFolio{579}{R}{MAGNETISCHE BEOBACHTUNGEN.}
\noindent nen, worin wir jedoch unsern Lesern nicht vorgreifen, dagegen aber die Erinnerung beifügen wollen, dass man bei allen erscheinenden Unähnlichkeiten vor allen Dingen die äussern Umstände sorgfältig erwägen muss, ehe man sie zur Grundlage von gewagten Vermuthungen macht. Als ein Beispiel kann die kleine Erhöhung dienen, die man in den graphischen Darstellungen des Augusttermins ---\quad---\quad---\quad--- für $18^{\mathrm u}5'$ bei den meisten Beobachtungsorten, am stärksten bei dem nordlichsten, Upsala, bemerkt. Dass dieselbe bei Dublin fehlt, oder nur eine schwache Spur davon sichtbar ist, ist allerdings merkwürdig, da kein Grund vorhanden ist, die Richtigkeit der Beobachtung selbst in Zweifel zu ziehen, und würde uns, zumal in Verbindung mit der vollständigen Erscheinung in Göttingen ---\quad--- zu interessanten Betrachtungen Anlass geben, wenn es überhaupt angemessen wäre, hier schon in solche uns einzulassen: allein dass diese Erhöhung auch bei Copenhagen fehlt, ist schlechterdings ohne alle Bedeutung, weil in diesem Termin in Copenhagen nur von 10 zu 10 Minuten, also um $18^{\mathrm u}5'$ gar nicht beobachtet ist.

Über die labyrinthischen Formen, welche die magnetischen Beobachtungen, bei Vereinigung der Declinations- und Intensitätsbewegungen in Einer Curve ---\quad---\quad---\quad---\quad--- uns vorführen, enthalten wir uns jeder Bemerkung hier nur deswegen, weil gegründete Hoffnung vorhanden ist, dass bald ein viel reicherer Stoff zu Gebote stehen wird. Wer inzwischen sich schon selbst in Betrachtungen über jene versuchen möchte, braucht sich wenigstens durch keine Zweifel an der Realität der durch das Bifilar-Magnetometer angezeigten Intensitätsbewegungen davon abhalten zu lassen. In der That sind solche Zweifel ganz unstatthaft geworden, nachdem bereits im Märztermin des gegenwärtigen Jahres 1838 ausser Göttingen noch an \textit{drei andern Orten} die gleichzeitigen Intensitätsbewegungen mit ähnlichen Bifilarapparaten beobachtet sind, und eine eben so bewundernswürdige Übereinstimmung gezeigt haben, wie wir seit vier Jahren an den Declinationsbewegungen zu finden gewohnt sind. Das Nähere darüber wird aber an die Bekanntmachung der Resultate der Beobachtungen von 1838 geknüpft bleiben müssen.
% Oxford 1867 copy, printed 579: 73\textsuperscript{*}; absent from the principal scan.


\clearpage
\SourceLeaf{586}\GaussVSourceFolio{580}{L}{MAGNETISCHE BEOBACHTUNGEN.}
\vspace*{.062\textheight}\pdfbookmark[1]{Astronomische Nachrichten Nr. 417. 1841. Februar 9.}{nachrichten-417}
\begin{center}
\vbox{\offinterlineskip
\centerline{\rule{.780\linewidth}{.35pt}}
\kern4pt
\centerline{\resizebox{.662\linewidth}{!}{\small \textsc{Schumacher}.\quad Astronomische Nachrichten Nr.\ 417.\quad 1841.\quad Februar 9.}}
\kern4.5pt
\centerline{\rule{.780\linewidth}{.35pt}}}
\end{center}

In einem öffentlichen Blatte fand ich unlängst die Nachricht, dass der Amerikanische Marine-Capitain \textsc{Wilkes}
\begin{center}\textit{dem magnetischen Südpole}\end{center}
\noindent ziemlich nahe gekommen sei, und dass er in $67^0\;4'$ südl.\ Breite und $147^0\;30'$ Länge (ohne Zweifel östlich von Greenwich) die magnetische Abweichung $12^0\;35'$ östlich, und die Neigung $87^0\;30'$ gefunden habe. Nach einer flüchtig angestellten Rechnung würde ich nun hienach einstweilen den wirklichen Pol
\begin{center}in $70^0\;21'$ südlicher Breite,\qquad $146^0\;17'$ Länge\end{center}
\noindent setzen. Dieser Platz liegt demjenigen, welchen meine Theorie [S.\ 163 d.\ B.] angegeben hat, viel näher, als ich selbst erwartet hatte. Der wirkliche Pol, wie ich dort vermuthet hatte, nordlicher, als der nach der Theorie berechnete; aber der Unterschied in der Breite erreicht nur den dritten Theil von dem, auf welchen ich nach Ansicht der Beobachtungen von Hobarttown gefasst war. Eben so liegt der wirkliche Pol westlicher, als der nach der Theorie berechnete, und hier ist der Unterschied fast genau so gross wie der a.\ a.\ O.\ von mir präsumirte. Übrigens ist unnöthig zu bemerken, dass in diesen hohen Breiten der Unterschied von sechs Längengraden nur eben so viel bedeutet, wie zwei Breitengrade.
\par\vspace{.07\textheight}\centerline{\rule{.255\linewidth}{.35pt}}

\clearpage
\SourceLeaf{587}\pdfbookmark[0]{Anzeigen nicht eigner Schriften.}{anzeigen-nicht-eigner}
\GaussVSourceUnnumbered
\vspace*{.26\textheight}
\begin{center}
{\LARGE\textls[140]{ANZEIGEN}}\par\vspace{.079\textheight}
{\large NICHT EIGNER}\par\vspace{.079\textheight}
{\LARGE\textls[480]{SCHRIFTEN.}}
\end{center}
\clearpage
\SourceLeaf{589}\GaussVSourceUnnumbered\fancyhead[L]{\small ANZEIGEN NICHT EIGNER SCHRIFTEN.}
\pdfbookmark[1]{Ueber die Daltonsche Theorie. 1830.}{daltonsche-theorie}
\vspace*{.079\textheight}
\begin{center}
\vbox{\offinterlineskip
\centerline{\rule{.600\linewidth}{.35pt}}
\kern4pt
\centerline{\resizebox{.507\linewidth}{!}{\small Göttingische gelehrte Anzeigen.\quad 1830 December 11.}}
\kern4.5pt
\centerline{\rule{.600\linewidth}{.35pt}}}
\end{center}

\textit{Ueber die \textup{\textsc{Dalton}}sche Theorie, von \textup{J. F.\ \textsc{Benzenberg}.} Düsseldorf, \textup{1830.} Bei J. E.\ Schaub.} (192 Seiten in 8., nebst drei Steindrucktafeln).

Die von \textsc{Dalton} aufgestellte Hypothese, dass die verschiedenen Gasarten, aus welchen die atmosphärische Luft besteht, gar nicht gegenseitig auf einander drücken, sondern eben so viele von einander gleichsam unabhängige Atmosphären ausmachen, hat bei wenigen Physikern bisher Beifall gefunden. Unter diesen zeichnet sich aber Herr \textsc{Benzenberg} durch den unermüdeten warmen Eifer, mit welchem er jene Hypothese seit beinahe zwanzig Jahren in Schutz nimmt, ganz besonders aus. Namentlich hat er in der \textsc{Daubuisson}schen trigonometrisch-barometrischen Messung der Höhe des Monte Gregorio einen wichtigen Grund \textit{für} die \textsc{Dalton}sche Hypothese gefunden. Es ist klar, dass die barometrischen Höhenmessungen, wenn die \textsc{Dalton}sche Hypothese wahr ist, anders berechnet werden müssen, als nach der gewöhnlichen Theorie. Bei dem 5260 Fuss hohen Monte Gregorio fand Herr \textsc{Benzenberg} das Resultat der ersten Rechnung um 16 Fuss kleiner, als nach der andern, und sehr nahe eben so viel übertraf letztere das Resultat der trigonometrischen Messung, welche Differenz mithin nach Herrn \textsc{Benzenberg}s Rechnung durch die Annahme der \textsc{Dalton}schen Hypothese fast vollkommen gehoben werden würde. Herr \textsc{Benzenberg} hat diese Rechnungen zuerst in \textsc{Gilbert}'s Annalen der Physik 1812 bekannt gemacht, und ist auch nachher
\par\clearpage
\SourceLeaf{590}\GaussVSourceFolio{584}{L}{ANZEIGEN.}
\noindent an andern Orten zu wiederholten Malen damit aufgetreten. Auch über andere Abschnitte der Physik, welche mit der \textsc{Dalton}schen Vorstellungsart in Berührung kommen, wie die Akustik und Eudiometrie, hat Herr \textsc{Benzenberg} sich ausgelassen, nicht sowohl, um Gründe \textit{für} jene Hypothese darin zu suchen, als vielmehr, um diejenigen Gründe, welche man daraus gegen dieselbe hernehmen kann und hergenommen hat, zu bekämpfen. In gegenwärtiger Schrift ist alles abermals vereinigt aufgestellt.

Die Anerkennung, welche diese Bemühungen bei den Physikern und Mathematikern gefunden haben, scheint, wenigstens was die ausländischen betrifft, nicht genügend gewesen zu sein. Herr \textsc{Benzenberg} erzählt in der Vorrede seines Buchs, dass er bei seinem Aufenthalt in Paris im Jahre 1815 \textsc{Gilberts} Annalen in der Bibliothek des Instituts gar nicht, und das Exemplar in der grossen Königlichen Bibliothek nicht aufgeschnitten gefunden habe. Mit \textsc{Laplace} habe er von seiner Theorie der Berechnung der Barometerhöhen in \textsc{Daltons} Hypothese gesprochen; allein dieser grosse Geometer sei damals schon alt gewesen.

Den deutschen Physikern kann doch der Vorwurf der Nichtbeachtung nicht mit Recht gemacht werden. Wir lesen in der neuen Ausgabe des physikalischen Wörterbuchs im Artikel Atmosphäre, dass `Herr \textsc{Benzenberg} der bedeutendste, gründlichste und eifrigste Vertheidiger der \textsc{Dalton}schen Theorie ist, dass mit Anwendung derselben die Höhen genauer, als ohne sie, berechnet werden, und dass darin ein bedeutender Beweis für die Richtigkeit derselben liege.'

Allein \textit{geprüft} scheint der Verfasser dieses Artikels die \textsc{Benzenberg}schen Rechnungen nicht zu haben: alle Zahlen sind nur ohne weiteres aus den \textsc{Gilbert}schen Annalen copirt. Dasselbe gilt von demjenigen, was über jene Rechnungen in dem Artikel Höhenmessung in dem erwähnten Wörterbuche gesagt ist. Vielleicht haben die Verfasser beider Artikel eine Prüfung deswegen für minder wesentlich gehalten, weil sie, den barometrischen Höhenmessungen überhaupt eine geringere Zuverlässigkeit beilegend, als Herr \textsc{Benzenberg}, die Beweiskraft von dessen Rechnungen doch nicht anerkannten, obwohl freilich der Verfasser des ersten Artikels dadurch das vorher angeführte zum Theil wieder aufhebt.

Es kann nicht die Absicht der gegenwärtigen Anzeige sein, unsere eigene Ansicht von der Zulässigkeit oder Unzulässigkeit der \textsc{Dalton}schen Hypothese selbst zu entwickeln, sondern wir beschränken uns auf dasjenige, was Herr~\textsc{Ben}-\par\clearpage
\SourceLeaf{591}\GaussVSourceFolio{585}{R}{BENZENBERG.\quad ÜBER DIE DALTONSCHE THEORIE.}
\noindent \textsc{zenberg} zur Unterstützung derselben in vorliegender Schrift von neuem vorgetragen hat, und namentlich auf seine Berechnung der barometrischen Höhenmessungen.

Es schien vor allen Dingen nöthig, erst die Richtigkeit dieser \textsc{Benzenberg}schen Berechnung selbst zu prüfen. Zu unserer Verwunderung ist daraus hervorgegangen, dass diese Berechnung unrichtig ist, und dass eine richtig geführte Rechnung ein ganz entgegengesetztes Resultat gibt.

Es wird nöthig sein, dies hier mit einiger Ausführlichkeit und so weit nachzuweisen, dass jeder in den Stand gesetzt wird, selbst zu urtheilen, um so mehr, da es sich hier nicht sowohl um einen Rechnungsfehler im eigentlichen Sinn, als um einen Irrthum im Räsonnement handelt.

Auf der zweiten und dritten Seite des Buchs finden wir zwei tabellarische Übersichten, jede von drei Columnen, welche wir einzeln durchgehen müssen (dieselben stehen in \textsc{Gilberts} Annalen der Physik, B.\ 42.\ S.\ 163 und 164).

In der ersten Columne wird, nach \textsc{Dalton}, die in 100 Theilen trockner atmosphärischer Luft enthaltene
\begin{center}
\begin{tabular}{@{}l@{\hspace{.119\linewidth}}c@{\hspace{.031\linewidth}}c@{\hspace{.031\linewidth}}c@{\hspace{.023\linewidth}}crl@{}}
\makebox[0pt][l]{Stickluft}&.&.&.&zu&$78.93$&Theilen\\
\makebox[0pt][l]{Sauerstoffluft}&&&.&--&$21.00$&\multicolumn{1}{c@{}}{--}\\
\makebox[0pt][l]{Kohlensaure Luft}&&&&--&$0.07$&\multicolumn{1}{c@{}}{--}
\end{tabular}
\end{center}
alles dem Raume nach, angegeben (Statt der zweiten Zahl steht im Buche selbst $21.90$, welches ein offenbarer Druckfehler ist).

Die zweite Columne setzt die specifischen Gewichte
\begin{center}
\begin{tabular}{@{}l@{\hspace{.191\linewidth}}c@{\hspace{.031\linewidth}}c@{\hspace{.031\linewidth}}c@{\hspace{.026\linewidth}}r@{}}
\makebox[0pt][l]{der Stickluft}&.&.&.&$0.9691$\\
\makebox[0pt][l]{der Sauerstoffluft}&&.&.&$1.1148$\\
\makebox[0pt][l]{der kohlensauren Luft}&&&&$1.5000$
\end{tabular}
\end{center}
das der gemeinen atmosphärischen Luft zur Einheit angenommen. Die erste und dritte Zahl sind von \textsc{Biot} entlehnt; die zweite hat Herr \textsc{Benzenberg}, wie er selbst sagt, so berechnet, `dass die Mischung genau das von \textsc{Biot} gegebene Gewicht trockner atmosphärischer Luft habe, nemlich $\frac{1}{10493}$ des Quecksilbers bei $0^0$ Reaumur und 28 Zoll Druck.' Man sieht, dass Herr \textsc{Benzenberg} sich selbst nicht klar gemacht hat, was er hier eigentlich hat rechnen wollen, denn das eben angezeigte ist \textit{dieser} Rechnung fremd, und offenbar kam es blos darauf an, der~Mi-\par\GaussVSourceGathering{74}
